\documentclass[10pt,a4paper]{article}
\usepackage{fontspec,amsmath,amssymb,graphicx}
\usepackage[margin=18mm]{geometry}
\setmainfont{Times New Roman}
\setlength{\parindent}{0pt}
\setlength{\parskip}{5pt}
\allowdisplaybreaks
\sloppy
\emergencystretch=2em
\begin{document}
\section*{BT\_6\_1\_1\_2\_03 — Lời giải bổ sung}
(i) $\sum\sin nx/n=(\pi-x)/2$ với $0<x<2\pi$, kéo dài chu kỳ $2\pi$; bằng 0 tại $x\in2\pi\mathbb Z$. Suy từ Fourier của hàm răng cưa hoặc phần ảo $-\log(1-re^{ix})$ rồi $r\uparrow1$.
(ii) Tích phân chuỗi (i), ấn định hằng bằng trung bình 0, cho $\sum\cos nx/n^2=\pi^2/6-\pi x/2+x^2/4$ trên $[0,2\pi]$, kéo dài tuần hoàn. Chuỗi hội tụ đều và tuyệt đối.
(iii) Tổng chuỗi hình học với $z=-e^{ix}/2$:
\[\sum_{n\ge1}(-1)^n2^{-n}\sin nx=\operatorname{Im}\frac z{1-z}=-\frac{2\sin x}{5+4\cos x}.\]
Chuỗi hội tụ tuyệt đối và đều.\end{document}